\subsection{Solving Absolute Value Equations} 
\begin{procedure}{Solving an Absolute Value Equation}
To solve an equation containing one absolute value, we use the following steps.
\begin{enumerate}
\item Use basic solving steps to isolate the absolute value.
\item Eliminate the absolute value by considering its properties.
\item Use basic solving steps to solve the resulting equation(s).
\end{enumerate}
\end{procedure}
\begin{Example}Solve each equation below.
\begin{lpicvsplit}{3}{9}
\foreach \x/\eqn in {0/2\Abs{x}+5=9,1/-\frac{2}{3}\Abs{5x-1}+7=3,2/2\Abs{\frac{3}{5}x-7}+3=1}
	\regtext{\x*\value{fcols}+\x+0.5*\value{fcols}}{-1}{$\eqn$};
\end{lpicvsplit}
\end{Example}
\begin{fact}
To eliminate the absolute value, once the equation is in the form $\Abs{X}=a$ for some real number $a$:
\begin{itemize}
\item If $a>0$, the equation has two solutions: $X=a$ or $X=-a$
\item If $a=0$, the equation has one solution: $X=0$
\item If $a<0$, the equation has no solution, since an absolute value cannot be negative.
\end{itemize}
\end{fact}